\chapter{Introduzione}
\label{chap:introduzione}

Descrizione del paradigma aggregate computing per sistemi distribuiti, le limitazioni
dovute ai sistemi embedded, cos'è zenoh e perchè è stato scelto (principalmente
parlando del modello peer-to-peer, non mi concentrerei troppo sulle performance
non avendo effettuato benchmark).

Descrizione di yaair, del modello a layer e del concetto di network.
\sidenote{come dovrei citare yaair essendo una libreria?}

\section{Obbiettivi}
\label{sec:obbiettivi}

Definizione chiara degli obbiettivi ed il focus di questo progetto:
\begin{itemize}
	\item Studio di Zenoh: il suo funzionamento, il modello peer
	      to peer (scouting, router, etc\dots) e come può essere utilizzato.
	\item Setup del techstack: embdedded rust, il framework esp-idf, integrazione
	      di librerie esterne e creazione del wrapper.
	      \sidenote{come dovrei usare inglesismi come techstack?}
	\item Creazione della network: comunicazione, gestione della perdita di nodi,
	      integrazione con yaair.
\end{itemize}

\section{Struttura della tesi}
\label{sec:struttura-della-tesi}

\note{Qui non so bene cosa mettere, è una sezione richiesta secondo il template
	mi verrebbe da riscrivere quello che ho già messo negli obbiettivi.}
